\documentclass[../Main.tex]{subfiles}
\begin{document}

\section{Specification of use case ``Whitelist Domain Management''}

\textbf{Use case name}: Whitelist Domain Management

\textbf{Actor}: Security Administrator

\textbf{Precondition}: The administrator has network access to the ShieldX-web dashboard.

\textbf{Flow of events}:
\begin{enumerate}
    \item The administrator navigates to the Whitelist Domains page via the sidebar.
    \item The system displays the current list of whitelisted domains in a table.
    \item To add a new domain:
    \begin{enumerate}
        \item The administrator clicks the ``Add Domain'' button.
        \item A modal dialog appears with input fields for domain name and optional notes.
        \item The administrator enters the domain (e.g., example.com) and clicks ``Add''.
        \item The system validates the domain name format and checks for duplicates.
        \item If valid, the domain is saved to the database and the table updates.
        \item If duplicate, an error toast notification is displayed.
    \end{enumerate}
    \item To edit an existing domain:
    \begin{enumerate}
        \item The administrator clicks the ``Edit'' button on the desired row.
        \item A modal dialog appears pre-filled with the domain's current data.
        \item The administrator modifies the fields and clicks ``Update''.
        \item The system validates and saves the changes.
    \end{enumerate}
    \item To delete a domain:
    \begin{enumerate}
        \item The administrator clicks the ``Delete'' button on the desired row.
        \item A confirmation dialog appears.
        \item The administrator confirms deletion.
        \item The domain is removed from the database and the table updates.
    \end{enumerate}
\end{enumerate}

\textbf{Postcondition}: The whitelist is updated in the database and propagated to agents during their next heartbeat cycle.

\section{Specification of use case ``Agent Heartbeat Monitoring''}

\textbf{Use case name}: Agent Heartbeat Monitoring

\textbf{Actor}: ShieldX Agent (primary), Security Administrator (secondary)

\textbf{Precondition}: The agent is running with network connectivity to the backend.

\textbf{Flow of events}:
\begin{enumerate}
    \item The agent collects system information (hostname, IP, MAC addresses, OS version).
    \item The agent sends a POST request to /heartbeat/ with its agent ID and system info.
    \item The backend checks if the agent ID exists in the database.
    \item If the agent exists, the backend updates last\_seen to the current timestamp and sets status to ``active''.
    \item If the agent is new, the backend creates a new Agent record with the provided information.
    \item The background scheduler periodically marks agents with last\_seen older than 5 minutes as ``offline''.
    \item The administrator can view the updated agent status on the dashboard.
\end{enumerate}

\textbf{Postcondition}: The agent's online status is tracked and visible to the administrator.

\section{Specification of use case ``Malware Alert Reporting''}

\textbf{Use case name}: Malware Alert Reporting

\textbf{Actor}: ShieldX Agent (primary), Security Administrator (secondary)

\textbf{Precondition}: The agent has detected a flow classified as malicious by the XGBoost model.

\textbf{Flow of events}:
\begin{enumerate}
    \item The agent captures network packets and assembles them into flows.
    \item For each completed flow, 28 statistical features are extracted.
    \item The XGBoost model classifies the flow as benign (0) or malicious (1).
    \item If the flow is classified as malicious:
    \begin{enumerate}
        \item The agent logs the detection locally with full flow details.
        \item The agent sends a POST request to /report-malware/ with agent ID, malware type, and flow features.
        \item The backend validates the agent ID and creates a MalwareAlert record.
        \item The backend returns the created alert with its assigned ID.
    \end{enumerate}
    \item If the flow is benign, the agent logs the flow as benign and continues monitoring.
    \item The administrator can view all alerts on the Malware Alerts page.
\end{enumerate}

\textbf{Alternative flow}: If the backend is unreachable, the agent continues local detection and retries alert transmission on the next heartbeat cycle.

\textbf{Postcondition}: Malicious flows are recorded in the database and visible to the administrator.

\end{document}
